% !TEX program = xelatex
% 2026 official-format smoke example.  Replace the prose with your verified
% results; do not add school, team, member, adviser, or logo information.
\def\GMCMContestEdition{二十三}
\documentclass[bwprint]{gmcmthesis}
\hypersetup{hidelinks}

\title{请替换为匿名论文题目}

\begin{document}
\maketitle

\begin{abstract}
本文档仅用于验证“华为杯”第二十三届中国研究生数学建模竞赛的匿名论文版式。正式摘要应简洁说明建模思路、主要方法、建立的模型、关键求解结果、结论与创新点，并只陈述能够由求解结果支撑的内容。摘要篇幅一般不超过两页，无需英文翻译。

\keywords{数学建模\quad 匿名提交\quad 格式核对\quad 可复现}
\end{abstract}

\section{问题重述}
此处用自己的语言概括题目任务、输入数据、约束条件和各小问的输出要求；不要复制身份信息或无关背景。

\section{模型假设}
模型假设应服务于问题并说明合理性。例如，先定义变量、集合和单位，再明确数据缺失、随机扰动或边界条件的处理方式。

\section{模型建立与求解}
在此给出可追溯的模型、推导、算法和结果。若引用公开资料，应在正文按出现次序使用方括号编号，并在参考文献中给出完整信息。

\section{模型评价与结论}
报告与问题对应的结果，说明验证、灵敏度或适用边界；不要用未验证的数字替代结果。

\begin{thebibliography}{9}
\bibitem{format}
“华为杯”第二十三届中国研究生数学建模竞赛论文格式规范，2026 年 9 月 16 日。
\end{thebibliography}

\end{document}
